%%%%%%%%%%%%%%%%%%%%%%%%%%%%%%%%%%%%%%%%%%%%%%%%%%%%%%%%%%%%%%%%%%%%%%%%%%%%%%%%%%%%%%%%%%%%%%%%
%% Corpora --- reframed draft, 28 November 2026.
%%
%% Organised by the new framing: the inclusion criterion is the FILE FORMAT, the screening is
%% auditable, and the unit of analysis is the TASK. Every structural fact below is traceable:
%%
%%   format contract, nine encodings .... docs/01-metodo-minimapas.md; reference/paper_02_encoding.tex
%%   screening, 27 candidates ........... reference/LEVANTAMENTO_DATASETS.md (13 Nov 2026),
%%                                       docs/00-origem-e-estado.md (20 Nov 2026 re-search),
%%                                       plus the 28 Nov 2026 pass recorded in docs/03-bases.md
%%   task slots per corpus .............. protocols/PROTOCOL-<cohort>.md, section "Tasks";
%%                                       counts cross-checked against the task_id column of
%%                                       runs/<cohort>/f1_within_descriptor/task_predictions.csv.gz
%%   participant and recording counts ... tables/cohorts.tex, generated from runs/
%%
%% TODO(before submission): the funnel of 27 -> 19 -> 8 -> 7 -> 6 and the task inventory of
%% Table~\ref{tab:tasks} should be emitted by scripts/x3_paper_tables.py like every other table in
%% this paper, so that the counts cannot drift from the runs. They are hand-set here only because
%% the generator does not yet have a task-level layer to read.
%%%%%%%%%%%%%%%%%%%%%%%%%%%%%%%%%%%%%%%%%%%%%%%%%%%%%%%%%%%%%%%%%%%%%%%%%%%%%%%%%%%%%%%%%%%%%%%%

\subsection{Corpora}
\label{sec:corpora}
%% Alias for the labels the shared generated tables reference. tables/*.tex are emitted by
%% scripts/x3_paper_tables.py for all three builds and their captions cite the earlier
%% draft's section names; aliasing here keeps one generator for three documents instead of
%% branching it per document.
\label{sec:cohorts}

\subsubsection{The inclusion criterion is the file format}

The method consumes a trajectory, not a dataset, so the corpora admitted to this study are defined
by what their files contain and not by what they are about.

\begin{definition}[Format contract]
\label{def:contract}
A corpus is admissible if it releases, for every recording and every sample $i$ of that recording,
the pen position $(x_i, y_i)$ and a timestamp $t_i$; and if it can be obtained by download, without
a request to an author. Pressure, azimuth, altitude and a pen-up/pen-down status are optional. Each
optional channel that is present adds encodings to the menu; each one that is absent removes them.
\end{definition}

Everything the pipeline needs follows from \cref{def:contract}. Speed is a finite difference of
$(x, y, t)$ and so is available wherever the contract holds. The nine encodings of the original
formulation \citep{prior2027minimaps} are the static trace, four single-signal minimaps --- speed,
pressure, azimuth, altitude --- and the four three-signal RGB triples that can be formed from those
signals, which requires the full channel set. A corpus with $x$, $y$ and $t$ alone therefore admits
two of the nine; one that adds a force channel admits three; one that adds pressure and a single
tablet-specific angle admits five; only the full set admits all nine. This is a property of the
release, not a defect of the corpus, and we report it per corpus rather than discarding the corpora
that constrain it: the narrow cases are the ones that test whether the encoding needs the whole
channel set to pay for itself.

The criterion is deliberately mechanical. ``Corpora resembling the one where the effect was
reported'' would let the published result choose its own test set, which is the failure mode this
study exists to avoid. The contract admits a smartphone screen, a child's dictation, a signature
and a single writer tracing twenty characters, all of which are in what follows, and none of which
we would have chosen for their resemblance to Parkinsonian handwriting.

\subsubsection{Screening, and what was rejected}
\label{sec:screening}

Twenty-seven candidate corpora were screened against \cref{def:contract}: an initial survey of the
handwriting-and-disease literature, two further independent searches, and a confirmatory pass
immediately before this draft. Nineteen fail the contract, for six distinct reasons, and we record
the reason for each because the reasons are themselves a finding about the state of this field:

\begin{itemize}
\itemsep0em
\item \textbf{Extracted attributes published instead of trajectories} (one corpus). The public
release gives per-task features, from which no image can be rendered.
\item \textbf{Static images of the product} (one). The release holds scans of the finished
handwriting, with no time axis at all.
\item \textbf{Supporting material that is not the signal} (three). The deposits hold acquisition
templates, analysis code, or aggregate measures; one Parkinson study's supporting information
consists of three figures and a methods document.
\item \textbf{Access conditional on an author request or on ethics approval} (seven). Several state
the intention to release and have not; one states explicitly that ethical and legal restrictions
prevent public sharing.
\item \textbf{Published address no longer serving the data} (one). The address given in the paper
redirects to a login screen.
\item \textbf{No public file of the raw series could be located} (six). The papers describe
per-sample acquisition; no corresponding download exists that we could find.
\end{itemize}

We stress what this list is not. It is not a judgement of those studies, and the request-to-author
route is legitimate; it is simply incompatible with a study whose point is that the test set was
not negotiated. Eight corpora were obtained, each as a published archive whose SHA-256 we recorded
on download, and the checksums are in the reproducibility package. Of the eight, seven satisfy
\cref{def:contract}; one does not, and we say so in \cref{sec:newhandpd} rather than quietly
dropping it. Of the seven, six carry a participant-level diagnostic or identity endpoint and enter
the task-level analysis; the seventh is a signature corpus whose label lives inside the writer and
is used, as \cref{sec:svc} explains, to verify that seven channels are read and nine encodings are
rendered correctly.

\subsubsection{The unit of analysis: thirty tasks}
\label{sec:tasks}

Each corpus elicits a fixed battery, and each participant contributes at most one recording per
battery slot. The slot --- the task --- is this study's unit of analysis, for the reasons given in
\cref{sec:intro}, and \cref{tab:tasks} is the complete inventory: thirty slots over six corpora.

The property the hypothesis turns on is assigned per slot, before any task-level quantity is
computed, by a rule that does not look at the data: a slot is \emph{figure imposed} when every
participant is asked to produce the same shape, so that what differs between the classes is the
execution, and \emph{shape informative} when the produced shape varies with what the participant
wrote, recalled or was asked to identify. Fifteen slots fall on each side. Two consequences of
applying the rule at slot level deserve to be stated rather than buried, because both move units
across the line relative to our own earlier corpus-level analysis. PaHaW is a mixed battery: its
Archimedean spiral is figure imposed, while its seven letter, word and sentence slots are shape
informative, so the corpus that was one ambiguous unit becomes one unit of the first kind and seven
of the second. Character Trajectories prescribes which character to write, but the character
\emph{is} the label, so all three of its slots are shape informative --- the extreme of that side,
and the reason the corpus is in the study.

% TODO(generator): emit this from scripts/x3_paper_tables.py. Counts verified on 28 Nov 2026
% against the task_id column of runs/<cohort>/f1_within_descriptor/task_predictions.csv.gz.
\begin{table}[htbp]
\caption{The thirty task-level units. \emph{Imposed} counts slots in which every participant
produces the same prescribed shape; \emph{informative} counts slots in which the produced shape
varies with the content. The assignment rule is fixed in \cref{sec:tasks} and applied before any
task-level metric. ParkinsonHW is listed on the sub-cohort its acquisition artefact leaves usable
(\cref{sec:floor}); its third acquisition block is excluded for the reason given in
\cref{sec:parkinsonhw}. Drot\'ar--Dobe\v{s} contributes a single slot because its released
recordings carry no task boundaries (\cref{sec:drotar}).}
\label{tab:tasks}
\begin{center}
\begin{small}
\begin{tabular}{llrrr}
\toprule
corpus & battery & slots & imposed & informative \\
\midrule
PaHaW & spiral, letters, words, sentence & 8 & 1 & 7 \\
ParkinsonHW & static and dynamic spiral & 2 & 2 & 0 \\
Zhu et al. & spiral, wave & 2 & 2 & 0 \\
DiaGraMo & 10 graphomotor, 4 writing/drawing & 14 & 10 & 4 \\
Drot\'ar--Dobe\v{s} & one unsegmented recording & 1 & 0 & 1 \\
Character Trajectories & three character pairs & 3 & 0 & 3 \\
\midrule
total & & \textbf{30} & \textbf{15} & \textbf{15} \\
\bottomrule
\end{tabular}
\end{small}
\end{center}
\end{table}

\begin{table}[t]
\caption{The cohorts analysed. \emph{Shape} is whether the produced figure carries the label, which \cref{sec:when} argues is what decides the outcome. \emph{Units} are participants and the parenthesis gives the class split, except in the two non-clinical cohorts, where the unit is the recording: Character Trajectories has a single writer, and in SVC2004 each of the $40$ writers contributes both classes, so a writer is not a class. $^{\dagger}$ParkinsonHW is reported on the only sub-cohort its own acquisition artefact leaves usable (36 patients from the second release plus the 15 controls, out of 76 participants in the full release); \cref{sec:floor} explains why.}
\label{tab:cohorts}
\vskip 0.1in
\begin{center}
\begin{small}
\begin{tabular}{llllrr}
\toprule
cohort & instrument & channels & shape & units & recordings \\
\midrule
ParkinsonHW$^{\dagger}$ & stylus, tablet & x, y, p, grip & fixed by task & 51 (15 / 36) & 98 \\
Zhu et al. & finger, phone & x, y only & fixed by task & 58 (30 / 28) & 116 \\
DiaGraMo & stylus, Cintiq & full & both & 273 (156 / 117) & 3796 \\
Drot\'ar--Dobe\v{s} & stylus, tablet & full & chosen by writer & 120 (57 / 63) & 120 \\
PaHaW & stylus, tablet & full & mixed & 75 (38 / 37) & 597 \\
Char.\ Trajectories & stylus, tablet & x, y, force & \textbf{is} the label & 853 (397 / 456) & 853 \\
SVC2004 & stylus, tablet & full & copied by forger & 1600 (800 / 800) & 1600 \\
\bottomrule
\end{tabular}
\end{small}
\end{center}
\vskip -0.1in
\end{table}


\subsubsection{PaHaW: the corpus the hypothesis comes from}
\label{sec:pahaw}

Seventy-five participants, thirty-seven with Parkinson's disease and thirty-eight controls, each
writing an eight-task battery on a digitising tablet that records all seven channels
\citep{drotar2016pahaw}. The battery progresses from drawing to writing: an Archimedean spiral
drawn from the inside out; the cursive letter \textit{l}, the bigram \textit{le} and the trigram
\textit{les}; the Czech words \textit{lektorka}, \textit{porovnat} and \textit{nepopadnout}, each
writable without lifting the pen; and the sentence \textit{Tramvaj dnes u\v{z} nepojede}, which
adds in-air transitions between words. Our copy holds 597 recordings: the spiral has 72 because
three are absent, and every other slot has one recording per participant. All nine encodings are
available. This is the corpus on which the encoding was reported \citep{prior2027minimaps}, so in
this study it is not evidence --- it is the regression control, included unchanged, and a change to
the pipeline that moves its numbers is a defect rather than an improvement.

\subsubsection{ParkinsonHW: the nearest replication, inside a contaminated release}
\label{sec:parkinsonhw}

Spiral drawing on a graphics tablet by Parkinson patients and controls, released on the UCI
repository \citep{isenkul2013parkinsonhw}. Two editions exist and they are the same collection:
thirty-five of the forty files of the earlier edition are byte-identical to files of the later one,
and the remaining five pair with named controls. We therefore treat the two editions as one cohort
of seventy-six participants, which a reader who takes them as independent would double-count.

The documented fields are $x$, $y$, $z$, pressure, a grip angle, a timestamp and a test
identifier. We discard $z$, which is zero in about ninety-five per cent of samples and undocumented,
and we do not treat the grip angle as the azimuth or the altitude of the original formulation,
because the repository's dictionary establishes no such correspondence; it enters as a third signal
under its own name. Five of the nine encodings are therefore available. The test identifier
separates three acquisition blocks, of which we keep two --- the static spiral, where the pattern
to trace is printed, and the dynamic spiral, where it blinks. The third is excluded, and the
decision was taken from the audit before any model was fitted: its pressure is identically zero for
whole blocks in twenty-three of the fifty-four participants who have it, which cannot describe a
pen pressing on a surface, and the readme describes it as circling a point rather than writing.

This corpus also carries the study's clearest acquisition artefact, and the analysis is reported on
the sub-cohort that artefact leaves usable --- the thirty-six patients of the second release plus
the fifteen controls, fifty-one participants and ninety-eight recordings. \Cref{sec:floor} gives
the measurement that forces the restriction; the point here is that the restriction was derived
from capture metadata alone, never from a classification score.

\subsubsection{Zhu et al.: a finger on a phone}
\label{sec:zhu}

Fifty-eight participants, twenty-eight with Parkinson's disease and thirty controls, drawing a
spiral and a wave with a bare fingertip on a smartphone screen \citep{zhu2025smartphone}, released
as supporting information to that paper: 116 trajectories with $x$ and $y$ in pixels and time in
milliseconds. There is no pressure and there are no pen angles, so exactly two of the nine
encodings exist here, the static trace and speed. Contact is implicit in a touch stream, so what
delimits a stroke had to be decided and written down before anything was rendered.

The scientific value of this corpus is precisely its poverty. It changes the instrument completely
--- no stylus, no tablet, no pressure --- while keeping the task property fixed, and it is the
cleanest test in the set of whether a speed channel alone, on two slots, carries anything.

\subsubsection{DiaGraMo: the only corpus that contains both kinds of task}
\label{sec:diagramo}

Two hundred and seventy-three Czech children, one hundred and fifty-six with handwriting
difficulties and one hundred and seventeen typically developing, on a Wacom Cintiq pen display that
records all seven channels \citep{zvoncakova2026diagramo}; 3,796 recordings over fourteen battery
slots after the release's own exclusion flag is applied. All nine encodings are available. Ten
slots impose the target figure --- three spirals, five loop tasks, a saw and a rainbow --- and four
do not: two Rey-complex-figure tasks, recognition and recall, then dictation and copy. The source
battery has sixteen tasks; the grade-3 and grade-4 versions of dictation and of copy are the same
task administered twice, so each pair shares a slot and no child contributes two recordings to one
slot.

This is the pivot of the study. Every other corpus varies the task property only by varying the
corpus, which varies the device and the population with it. Here the same children, the same
display and the same session produce both kinds of slot, so the contrast the hypothesis is about
can be formed with everything else held fixed. It is also the corpus in which our own earlier,
corpus-level version of that contrast failed, with the margin larger on the side where the
hypothesis predicted none (\cref{sec:results-diagramo}) --- which is the main reason the present study moves
to the slot as its unit. The group that collected the material has analysed it with
direction-based fractional-order derivative features \citep{gavenciak2025fractional}; we cite that
work to place the corpus and compare against pipelines we ran ourselves rather than against its
reported numbers.

\subsubsection{Drot\'ar--Dobe\v{s}: the fullest channel set and the coarsest unit}
\label{sec:drotar}

One hundred and twenty children, fifty-seven with dysgraphia and sixty-three controls, recorded on
a tablet with all seven channels and explicit pen-up/pen-down status
\citep{drotar2020dysgraphia}. All nine encodings are available, and this is the second independent
dysgraphia cohort in the study, which matters because a single dysgraphia result would be a
property of one clinic.

One recording per child spans several elicited tasks --- repeated letters, syllables, words and a
sentence --- and the released files carry no task boundaries. Each recording is therefore treated
as one whole and contributes a single slot, which is why this corpus is one row of \cref{tab:tasks}
rather than five. The consequence is accepted rather than worked around: no per-task comparison is
made here, and locating those boundaries is named in \cref{sec:limitations} as the obvious next
study rather than attempted mid-analysis, since a segmentation tuned after seeing a score would
compromise everything else.

\subsubsection{Character Trajectories: where the shape is the label}
\label{sec:chartraj}

Two thousand eight hundred and fifty-eight trajectories of twenty characters by a single writer on
a tablet that also records tip force \citep{williams2006chartraj}. Three encodings are available:
the static trace, speed, and force in the place of pressure.

Two properties make it the deliberate null end of the hypothesis. First, the label is the shape: a
reader identifies a character from ink alone, so there is by construction nothing for a painted
channel to add that the geometry does not already carry. Three character pairs chosen for shape
similarity --- \textit{u}/\textit{v}, \textit{o}/\textit{a}, \textit{m}/\textit{n} --- were declared
before any model was fitted and are the three slots of \cref{tab:tasks}, 853 recordings in all; the
canonical twenty-class problem over all 2,858 recordings is reported as well, registered as an
addendum to that corpus's protocol before it ran. Second, there is one writer, so the unit is the
recording, the participant-level fusion is the identity, and no participant-level claim is made
here. We report that as a structural difference rather than smoothing over it.

What the release contains is also unusual, and it inverts the contract: the stored quantity is
velocity, already differentiated, Gaussian-smoothed and normalised by the authors, with the
trajectory recovered by integration at our end. This is the one corpus in the study whose trace is
a reconstruction and whose dynamics are the primary record. Its timestamps are likewise
reconstructed at a fixed rate and it has no clock origin, which makes its acquisition floor chance
by construction --- and therefore makes it the verification of the floor procedure itself
(\cref{sec:floor}).

\subsubsection{SVC2004: the rendering and reading check}
\label{sec:svc}

One thousand six hundred signatures from forty writers, with the same seven fields as PaHaW
\citep{yeung2004svc}. All nine encodings render, which is what the corpus is for: if the seven
channels are read in the wrong order, in the wrong units or with the wrong angle convention, this
is where it shows, because the file specification is published and unambiguous. Its label, genuine
against skilled forgery, lives inside the writer rather than identifying one, so it has no
participant endpoint and no acquisition floor, and it is excluded from the task-level analysis by
construction. It is also a conceptually interesting positive control --- a skilled forger copies
the shape deliberately, so only the dynamics remain --- and we report it as such
(\cref{sec:results-svc}) without letting it into the estimand.

\subsubsection{NewHandPD: obtained, and rejected by the contract}
\label{sec:newhandpd}

Seven hundred and ninety-two signals from sixty-six participants, thirty-one with Parkinson's
disease and thirty-five controls, recorded by a sensor-equipped pen at a kilohertz. We obtained it, and it
fails \cref{def:contract}: the release gives six channels of pen-mounted sensor output and no
position of the tip on the writing surface. There is no trajectory, so there is no image, and no
result of this study can be read as evidence about this corpus. We report the rejection because
this corpus is frequently listed alongside the others as a handwriting benchmark, and the
distinction between ``dynamic handwriting data'' and ``a sampled trajectory'' is exactly the
distinction \cref{def:contract} is drawing.

\subsubsection{Availability, licences and what we may redistribute}
\label{sec:availability}

Every corpus analysed here is public, and none of it is ours. We redistribute no recording.
The reproducibility package releases the code, the per-corpus protocol documents, the configuration
of every run, the aggregate result tables and the SHA-256 of each archive as downloaded, which is
what allows the pipeline to be re-run by a reader who obtains the archives themselves. Per-participant
predictions are not published. SVC2004 is distributed for research use with commercial use
excluded; the remaining archives carry their own repository or journal terms, which we follow.
